\chapter{Conclusiones y Trabajo Futuro}
\label{ch:conclusiones}

Este capítulo sintetiza los principales hallazgos de la investigación, estructurados en base a los factores de variación analizados. Finalmente, se plantean diversas líneas de trabajo futuro derivadas de los resultados obtenidos.

\section{Conclusiones}

\subsection{Algoritmos}

\begin{itemize}
    \item \textbf{NSGA-II} exhibe un alto rendimiento general en el experimento absoluto (dominando en solitario las métricas de HV, Rango y SumaMín, y coliderando el primer estrato de significancia en IGD+, MínSuma y Hmed) impulsado por su operador de distancia de agrupamiento, que demuestra un funcionamiento óptimo en este espacio discreto y altamente correlacionado. Sin embargo, bajo el modelo proporcional, su rendimiento disminuye en prácticamente todas las métricas, cediendo en muchos casos el primer puesto a AGE-MOEA-II (cayendo a estratos estadísticos inferiores en MínSuma, Hmed y GD+, y perdiendo su dominio en solitario en HV y SumaMín). Como se ha demostrado empíricamente, este decrecimiento se debe a la descorrelación de los objetivos inducida por la formulación proporcional, la cual aumenta la dimensionalidad geométrica del frente y reduce la eficacia de su distancia de agrupamiento.
    \item \textbf{AGE-MOEA-II} es el diseño más beneficiado por la transición proporcional. Evoluciona de un rendimiento secundario en el escenario absoluto a coliderar el primer nivel estadístico en la mayoría de métricas clave (HV, IGD+, SumaMín y Tmáx) y a dominarlo en solitario en MínSuma y Hmed. Su capacidad intrínseca para estimar la curvatura del frente, unida a la inyección constante de soluciones dominadas, le confiere la flexibilidad geométrica necesaria para adaptarse a la mayor complejidad y falta de correlación del experimento proporcional.
    \item \textbf{Los MaOEA (NSGA-III y U-NSGA-III)} se asientan en un estrato intermedio de rendimiento para ambos experimentos, mostrando únicamente ligeras mejoras de densidad poblacional bajo el experimento proporcional. Su dependencia de vectores de referencia, pese a estar diseñados para MaOPs, no contribuye a lograr un alto rendimiento. Son dos algoritmos muy similares que tienen resultados muy parecidos.
    \item \textbf{SPEA2} ofrece resultados intermedios, algo inferiores a los de los MaOEA (NSGA-III y U-NSGA-III), posicionándose por detrás de ellos en algunas de las métricas (MínSuma y GD+ para ambos experimentos, en SumaMín absoluto y en Hmed proporcional), equivaliéndose estadísticamente en otras (como en SumaMín proporcional), y superándolos de forma puntual (liderando la diversidad Hmed absoluto).
    \item \textbf{SMS-EMOA} reproduce un patrón estable en ambos experimentos: se asienta en el primer estrato estadístico de las métricas de convergencia (1ª posición para GD+) y de tolerancia biológica (Tmed y Tmáx), pero queda relegado a las últimas posiciones (generalmente el penúltimo estrato) en métricas integrales, de diversidad y balance (HV, IGD+, Rango, MínSuma y SumaMín). Su estricta arquitectura de evolución de estado estacionario ($\mu+1$) desencadena una pérdida de diversidad genética, provocando un estancamiento prematuro documentado en ambos modelos.
    \item \textbf{RVEA} se estanca en el último nivel de significancia estadística en ambos experimentos según la mayoría de las métricas. Su limitación de diseño (el uso de la métrica de distancia de penalización angular sin mecanismos de normalización dinámica) provoca que gran parte de los individuos de su población se hallen en regiones infactibles. Aunque la compresión del experimento proporcional mitiga levemente este problema, su desempeño general sigue siendo el más bajo del conjunto.
    \item Atendiendo a la magnitud de su influencia estadística, el algoritmo explica en promedio un $16{,}67\%$ de la varianza del rendimiento, actuando como un optimizador de segundo nivel, por detrás de la estrategia de inicialización poblacional.
\end{itemize}

\subsection{Inicializaciones}

\begin{itemize}
    \item \textbf{\textit{Greedy Ting}:} Domina HV proporcional, IGD+, Rango y SumaMín. Su formulación híbrida —que combina soluciones voraces con aleatorias— es una inicialización que tiene un rendimiento similar a Random Dense, aunque esta proporciona al algoritmo una mayor diversidad y para el experimento proporcional ofrece mejores resultados.
    \item \textbf{\textit{Random Dense}:} Lidera HV absoluto, MínSuma, Tmed y Hmed. Esta estrategia de inicialización obtiene mejores resultados en el experimento absoluto, y tiene un rendimiento similar a la anterior.
    \item \textbf{\textit{Greedy Multi}:} Queda relegada al último puesto en rendimiento general, logrando destacar de forma aislada en GD+ (convergencia puntual) y Tmáx (tolerancia máxima). Dado que ambas métricas premian la búsqueda unidimensional hacia paneles muy pequeños, esta inicialización demuestra carecer de la amplitud necesaria para guiar la optimización global del frente de Pareto.
    \item En síntesis, la experimentación concluye que no existe una inicialización universalmente superior. Es un factor que depende de la formulación de las funciones objetivo y de las preferencias del usuario.
    \item La condición de parámetro primario queda demostrada al medir su tamaño de efecto: la inicialización es responsable del $64{,}46\%$ de la variabilidad total de los resultados, siendo cuatro veces más determinante que el propio algoritmo.
\end{itemize}

\subsection{Operadores de cruce}

\begin{itemize}
    \item A lo largo de los dos diseños experimentales y en la inmensa mayoría de métricas analizadas (salvo el caso puntual de la métrica Hmed proporcional, donde la inferencia estadística no determina variaciones), UX y HUX se asientan monolíticamente en el estrato de máximo rendimiento estadístico, haciendo que ambos operadores sean intercambiables en la práctica.
    \item El operador de cruce de dos puntos (2P) se clasifica consistentemente en un segundo nivel de significancia, seguido del cruce de un punto (1P) en el estrato inferior.
    \item Si bien es el componente estructural con menor grado de impacto global (apenas un $4{,}51\%$ de la varianza total del rendimiento), la selección del operador de cruce tiene un efecto estadísticamente significativo.
\end{itemize}


\subsection{Impacto de la formulación de funciones objetivo (Absoluto vs. Proporcional)}

\begin{itemize}
    \item \textbf{Descorrelación matemática y dimensionalidad:} El análisis estadístico de Spearman demuestra que la formulación absoluta genera un espacio de objetivos fuertemente redundante y sinérgico ($\rho \geq 0{,}96$), comportándose en la práctica como un problema de baja dimensionalidad. Por el contrario, la división proporcional introducida originalmente por Ting \textit{et al.} fractura esta dependencia lineal, induciendo conflictos reales (con magnitudes de correlación que caen hasta $|\rho| = 0{,}57$) e incrementando drásticamente la dimensionalidad efectiva del frente de Pareto. Esta mutación geométrica justifica la elección matemática de la formulación y es la responsable directa de alterar el comportamiento de los algoritmos.
    \item \textbf{Beneficiados:} En la mayoría de las métricas, AGE-MOEA-II asciende al primer estrato estadístico al adaptarse orgánicamente a la nueva dimensionalidad; los MaOEAs (NSGA-III y U-NSGA-III) ganan cardinalidad poblacional en sus frentes; y RVEA mitiga parcialmente su desplazamiento hacia soluciones infactibles.
    \item \textbf{Perjudicados:} NSGA-II pierde su liderazgo del experimento absoluto a causa de la ineficacia de la distancia de agrupamiento en espacios descorrelacionados, compartiendo el primer puesto con AGE-MOEA-II o cediéndolo según la métrica.
    \item La compacidad media de los frentes de Pareto se reduce en todas las estrategias de inicialización (descendiendo de $\sim 345$ a $\sim 269$ SNPs para \textit{Random Dense}, de $\sim 232$ a $\sim 118$ SNPs para \textit{Greedy Ting}), y de $\sim 30$ a $\sim 22$ SNPs para Greedy Multi, lo que demuestra que la formulación proporcional penaliza los paneles de gran tamaño.
    \item Los ranquines de rendimiento para inicializaciones y operadores de cruce se mantienen mayoritariamente estables entre ambas formulaciones, si bien es cierto que es apreciable que Greedy Ting rinde mejor en el experimento proporcional que en el absoluto, obteniendo el primer puesto en HV proporcional.
\end{itemize}

\subsection{Otros hallazgos}

\begin{itemize}
    \item \textbf{La saturación prematura de la dominancia de Pareto:} Con la excepción de AGE-MOEA-II, los algoritmos regidos por mecanismos de ordenación por dominancia sufren un agotamiento exploratorio temprano frente a los cuatro objetivos del problema. De media, NSGA-II clasifica al 100\% de sus individuos en el primer frente no dominado en tan solo 36 generaciones, mientras que SPEA2 lo hace en 45 y SMS-EMOA en apenas $11{,}85$. A partir de ese umbral, los algoritmos abandonan la evaluación por rango de Pareto y delegan la evolución a sus operadores secundarios de densidad espacial (distancia de agrupamiento, \textit{k-NN} y aportación marginal al HV respectivamente).
    
    Por el contrario, AGE-MOEA-II destaca como el único del estudio capaz de dividir su población en frentes de Pareto de diferente ranquin durante todas las generaciones. Esto garantiza el mantenimiento permanente de una reserva genética de soluciones dominadas. Esta inyección continuada de diversidad evita el estancamiento y es en parte responsable de su éxito.
    \item \textbf{Estabilización bajo el criterio de parada:} Las ejecuciones de 500 generaciones son suficientes para abarcar el proceso evolutivo de los algoritmos más robustos en este dominio. Las variaciones marginales registradas en el hipervolumen a lo largo del último quinto de la simulación (generaciones 400 a 500) se cuantifican en un incremento medio de apenas un $+1{,}06\%$ para la totalidad de configuraciones, demostrando la convergencia práctica de los frentes resultantes y desaconsejando incurrir en costes computacionales adicionales.
    \item \textbf{Reproducibilidad y comparativa con el estado del arte:} El contraste de los resultados obtenidos con la literatura previa (Moqa \textit{et al.}) revela discrepancias significativas tanto en el rendimiento de los algoritmos como en el impacto de la inicialización. Sin embargo, se concluye que resulta inviable determinar la causa exacta de estas divergencias debido a la indisponibilidad del \textit{software} original del estudio de referencia. Únicamente mediante el acceso al código fuente, donde residen los detalles importantes de implementación no documentados (tales como la lógica interna de la inicialización voraz, la gestión de individuos infactibles o los métodos de normalización aplicados), se podrían aislar de forma rigurosa los factores que originan estas discrepancias.
\end{itemize}

\section{Trabajo Futuro}

A partir de las limitaciones identificadas durante el presente trabajo de investigación, se plantean diversas líneas de continuidad enfocadas a perfeccionar y extender la aplicabilidad de los EAs sobre el problema genético de selección de etiquetas SNP:

\begin{itemize}
    \item \textbf{Estudio en profundidad de estrategias de inicialización:} Los resultados de este trabajo demuestran que la inicialización es el factor con mayor impacto sobre el rendimiento final. Esta evidencia motiva el diseño y evaluación de nuevas estrategias de inicialización poblacional adaptadas a este problema, explorando heurísticas basadas en información biológica previa del conjunto de datos (como frecuencias alélicas o patrones de desequilibrio de ligamiento) que puedan proporcionar puntos de partida más útiles que las inicializaciones evaluadas en este estudio, que no analizan previamente el conjunto de datos.
    \item \textbf{Mejora de las limitaciones algorítmicas identificadas:} Para los algoritmos basados en descomposición, especialmente RVEA, cuyo rendimiento ha sido el más bajo del estudio debido a su penalización angular sin normalización dinámica, se propone investigar mecanismos de adaptación de los vectores de referencia que mitiguen el desplazamiento hacia regiones infactibles. Por otro lado, SMS-EMOA requiere operadores de preservación de diversidad que frenen su pérdida de material genético. Asimismo, se plantea como línea de trabajo la implementación en pymoo de una versión de MOEA/D compatible con el manejo de restricciones, algoritmo que ha quedado excluido de este estudio por incompatibilidad con la restricción de tolerancia del problema.
    \item \textbf{Parametrización fina:} Llevar a cabo estudios de sensibilidad dirigidos a los factores estructurales de los algoritmos de alto rendimiento. Se propone evaluar de forma integrada el impacto de la probabilidad de cruce (UX/HUX) y de mutación, el tamaño de la población ($N$), el tamaño del torneo, y el diseño de tasas de mutación dinámicas destinadas a mitigar el estancamiento evolutivo en fases tardías.
\end{itemize}